\section{符号说明}
\begin{table}[htbp]
\centering
\caption{主要符号}
\begin{tabularx}{0.95\textwidth}{>{\raggedright\arraybackslash}p{2.2cm}X}
\toprule
符号 & 含义 \\
\midrule
$R_j(t)$ & 第 $j$ 类底层资源：水草、浮游植物、浮游动物、底栖饵料 \\
$C_j(t)$ & 第 $j$ 类四大家鱼功能群：草鱼、鲢鱼、鳙鱼、青鱼 \\
$M(t)$ & 中层共享猎物代理：中小型鱼类、幼鱼与鱼虾饵料 \\
\makecell[l]{$D(t),A(t),$\\$V(t)$} & 分别表示长江江豚、长江鲟和外来入侵强度 \\
\makecell[l]{$B_{\text{carp}}(t),$\\$R_{\text{basal}}(t)$} & 分别表示四大家鱼总量与底层资源总量 \\
$\mathrm{CPUE}^*(t)$ & 四大家鱼单网次捕获量代理（权重 1,1,1,0.7） \\
\makecell[l]{$F,u,$\\$F_0$} & 分别表示捕捞死亡率、污染强度和反事实捕捞压力 \\
$\delta_B,\gamma_B$ & 通道阻隔损失及长江鲟敏感性放大因子 \\
\makecell[l]{$k_M,R_A,$\\$K$} & 分别表示中层猎物承载力基准、人工放流强度和三级食物链承载力 \\
$\Phi(t)$ & 食压指数，刻画消费者对底层资源的相对压力 \\
$\lambda_{\max}$ & 最大李雅普诺夫指数 \\
$E(t)$ & 综合生态状态指数 \\
\bottomrule
\end{tabularx}
\end{table}

